\newpage

\begin{center}
  {\large\bfseries\color{titlegreen} এক নজরে মুনশাইবের বাব}

  \vspace{0.12cm}

  {\arabicfont\normalsize\bfseries\textarabic{ثلاثة وأربعون بابا}}
\end{center}

\vspace{0.15cm}

\noindent রূপান্তরশীল ফেল মূল অক্ষরের হিসাবে দুই ভাগ:
ছুলাসী (মাজির মূল অক্ষর তিন) ও রুবায়ী (মূল অক্ষর চার)।

\vspace{0.25cm}

\begingroup
\setlength{\arrayrulewidth}{0.5pt}
\setlength{\tabcolsep}{3pt}
\renewcommand{\arraystretch}{1.45}
\small

\noindent\begin{tabularx}{\textwidth}{|C|C|}
  \hline
  \rowcolor{headergreen}
  {\color{white}\textbf{নাম}} &
  {\color{white}\textbf{মানে}} \\
  \hline
  \rowcolor{lightgreen}
  মুজারাদ & মূলের সঙ্গে বাড়তি অক্ষর নেই \\
  \hline
  মাযীদ ফিহি & মূলের সঙ্গে বাড়তি অক্ষর আছে \\
  \hline
  \rowcolor{lightgreen}
  মুলহাক & ছুলাসী মাযীদ, চলে রুবায়ীর ওজনে \\
  \hline
  মুতাররাদ & ওজন বেশি চলে \\
  \hline
  \rowcolor{lightgreen}
  শায & ওজন কম চলে \\
  \hline
\end{tabularx}

\vspace{0.35cm}

\noindent\begin{tabularx}{\textwidth}{|C|C|}
  \hline
  \rowcolor{headergreen}
  {\color{white}\textbf{দল}} &
  {\color{white}\textbf{সংখ্যা}} \\
  \hline
  \rowcolor{lightgreen}
  ছুলাসী মুজারাদ মুতাররাদ & ৫ \\
  \hline
  ছুলাসী মুজারাদ শায & ৩ \\
  \hline
  \rowcolor{lightgreen}
  ছুলাসী মাযীদ, হামযায় ওয়াসল & ৯ \\
  \hline
  ছুলাসী মাযীদ, হামযায় কাতঈ ও তা-আলিফ & ৫ \\
  \hline
  \rowcolor{lightgreen}
  রুবায়ী মুজারাদ & ১ \\
  \hline
  রুবায়ী মুনশাইব & ৩ \\
  \hline
  \rowcolor{lightgreen}
  মুলহাক বি রুবায়ী মুজারাদ & ৭ \\
  \hline
  মুলহাক বি তাদাহরুজ & ৮ \\
  \hline
  \rowcolor{lightgreen}
  মুলহাক বি ইহরানজাম & ২ \\
  \hline
  \rowcolor{headergreen}
  {\color{white}\textbf{মোট}} &
  {\color{white}\textbf{৪৩}} \\
  \hline
\end{tabularx}
\endgroup

\vspace{0.4cm}

\begin{center}
  {\bfseries ছুলাসী মুজারাদ}
  \enspace
  {\arabicfont \textarabic{ثلاثي مجرد}}
\end{center}

\vspace{0.1cm}

\noindent চিহ্ন আইনের হরকত। মাজি ও মুজারে দুই হরকত মিলিয়ে বাব চেনা যায়।

\vspace{0.2cm}

\begingroup
\setlength{\arrayrulewidth}{0.5pt}
\setlength{\tabcolsep}{3pt}
\renewcommand{\arraystretch}{1.55}
\small

\noindent\begin{tabularx}{\textwidth}{%
    |>{\hsize=0.28\hsize\centering\arraybackslash}X%
    |>{\hsize=1.22\hsize\centering\arraybackslash}X%
    |>{\hsize=1.50\hsize\raggedright\arraybackslash}X|}
  \hline
  \rowcolor{headergreen}
  \multicolumn{3}{|c|}{\color{white}\textbf{মুতাররাদ}} \\
  \hline
  \rowcolor{headergreen}
  {\color{white}\textbf{বাব}} &
  {\color{white}\textbf{মাজি · মুজারে}} &
  {\color{white}\textbf{চিহ্ন}} \\
  \hline
  \brow{lightgreen}{১}{%
    {\arabicfont \textarabic{نَصَرَ يَنْصُرُ}}\newline সাহায্য করা}{%
    মাজি যবর, মুজারে পেশ}
  \brow{white}{২}{%
    {\arabicfont \textarabic{ضَرَبَ يَضْرِبُ}}\newline প্রহার করা}{%
    মাজি যবর, মুজারে যের}
  \brow{lightgreen}{৩}{%
    {\arabicfont \textarabic{سَمِعَ يَسْمَعُ}}\newline শোনা}{%
    মাজি যের, মুজারে যবর}
  \brow{white}{৪}{%
    {\arabicfont \textarabic{فَتَحَ يَفْتَحُ}}\newline খোলা}{%
    মাজি ও মুজারে যবর}
  \brow{lightgreen}{৫}{%
    {\arabicfont \textarabic{كَرُمَ يَكْرُمُ}}\newline সম্মানিত হওয়া}{%
    মাজি ও মুজারে পেশ। লাযিম; ইসমে ফায়েল প্রায়
    {\arabicfont \textarabic{فَعِيلٌ}} ওজনে}
  \rowcolor{headergreen}
  \multicolumn{3}{|c|}{\color{white}\textbf{শায}} \\
  \hline
  \brow{lightgreen}{১}{%
    {\arabicfont \textarabic{حَسِبَ يَحْسِبُ}}\newline ধারণা করা}{%
    মাজি ও মুজারে যের}
  \brow{white}{২}{%
    {\arabicfont \textarabic{فَضِلَ يَفْضُلُ}}\newline উপস্থিত হওয়া}{%
    মাজি যের, মুজারে পেশ। সহীহ মাসদার আসে না}
  \brow{lightgreen}{৩}{%
    {\arabicfont \textarabic{كَادَ يَكَادُ}}\newline কাছে আসা}{%
    মাজি পেশ, মুজারে যবর। মধ্যবর্ণ ইল্লাত}
\end{tabularx}
\endgroup
